% Part: first-order-logic
% Chapter: axiomatic-deduction
% Section: proving-things-quant

\documentclass[../../../include/open-logic-section]{subfiles}

\begin{document}

\olfileid{fol}{axd}{prq}

\olsection{带量词的推导}

\begin{ex}
我们来给出 $(\lforall[x][!A(x)] \land
\lforall[y][!B(y)]) \lif \lforall[x][(!A(x) \land !B(x))]$ 的一个推导。

首先，注意
\begin{align*}
  (\lforall[x][!A(x)] \land \lforall[y][!B(y)]) & \lif \lforall[x][!A(x)]\\
  \intertext{是 \olref[prp]{ax:land1} 的一个实例，而}
  \lforall[x][!A(x)] & \lif !A(a) \\
  \intertext{是 \olref[qua]{ax:q1} 的一个实例。因此，根据 \olref[pro]{prop:chain}，我们知道}
  (\lforall[x][!A(x)] \land \lforall[y][!B(y)]) & \lif !A(a) \\
  \intertext{是可推导的。类似地，由于}
  (\lforall[x][!A(x)] \land \lforall[y][!B(y)]) & \lif \lforall[y][!B(y)] \qquad\text{和}\\
    \lforall[y][!B(y)] & \lif !B(a)\\
    \intertext{分别是 \olref[prp]{ax:land2} 和 \olref[qua]{ax:q1} 的实例，}
    (\lforall[x][!A(x)] \land \lforall[y][!B(y)]) & \lif !B(a)\\
    \intertext{根据 \olref[pro]{prop:chain} 也是可推导的。使用 \olref[prp]{ax:land3} 的一个适当实例，并应用两次~\MP，可知}
    (\lforall[x][!A(x)] \land \lforall[y][!B(y)]) & \lif (!A(a) \land !B(a))\\
    \intertext{是可推导的。我们现在可以应用量词规则 \QR{} 来得到}
(\lforall[x][!A(x)] \land \lforall[y][!B(y)]) & \lif \lforall[x][(!A(x) \land !B(x))].
\end{align*}
\end{ex}

\end{document}
